% EGA V, English translation by P. Blass, 1987 (cote 24): archive cote 24, site page 5 (PDF page 6).
% typescript. Draft compared: none.
% AI-transcribed by Claude Opus 5.5 (claude-opus-5-5, adaptive thinking, effort high), Message Batches API, batch ega24b (msgbatch_01RGJJeHdR96U16ZWFuUjtDx), 2026-11-13; not yet checked by eye.
\scanpage{24/5}
\mspage{5}

(i.e. $\underline{L}$ or $\underline{L}$-1 \hand{\unsure{$\subset$|$\in$}} $\underline{E}$) and $\phi$ is obtained by taking for $\phi$ a
section of $\underline{L}$-1 which does not vanish at any point i.e. by a
trivialization of $\underline{L}$-1 (which exists in every case locally). Let us
note that H$\xi$ is nothing else but P($\underline{E}$ / $\underline{L}$-1) (canonical isomorphism)
that is a third way of describing H$\xi$ (N.B. P($\underline{E}$ / $\underline{L}$-1) is indeed
canonically embedded in P = P($\underline{E}$) which has the advantage of
proving in addition that H is a projective fibration over S and is
a fortiori smooth over S. (Again it would be necessary to say in
par 17 that a projective fibration is smooth). Without a doubt it
would be better to begin with that one. [Ask AG par 17 of what?]

\uline{Remarks}

The formation of H$\xi$ from $\xi$ is compatible with base change in
other words one finds a homomorphism of functors [Ask AG about
next symbol]

\quad $^{\circ}$

(Sch) /S $\rightarrow$ (Ens), PV $\rightarrow$ \uline{Div}(P/S) where the second term denotes the
functor of ``relative divisors'' of P/S whose values at S' (an
arbitrary S prescheme) is the set of closed subschemes of $\underline{P}_{S}$' which
are complete intersections transversal and of codimension 1
relative to S' (compare Par. 19).

It is easy to show that this homomorphism of functors is a
monomorphism, in other words that $\xi$ is determined of H$\xi$. (This
last fact justifies the terminology ``scheme of hyperplanes'' used
above.) We shall see that the functor \uline{Div}(P/S) is representable
by the prescheme (direct) sum of P(Symm$^{k}$(E)) so that $\underline{P}$ can be
identified to an open and closed subscheme of \uline{Div}(P/S).....
